% GENERATED FILE. Edit canonical lessons, then run npm run generate:books.
% canonical-source-hash: fnv1a64:c5099756483c9352
% canonical-lessons: BN-C327-uponyas, BN-C327-jiboni, BN-C327-grontho, BN-C327-odhyay, BN-C327-prokashok

\chapter{Books to Remember}
\label{ch:bn-books-to-remember}
\hlchaptermodality{327}

\begin{chapteropening}
\textbf{By the end of this chapter:} \emph{I can talk about novels, biographies and the people who publish them.}

Complete the last lesson of chapter 327: I can talk about novels, biographies and the people who publish them.
\end{chapteropening}

\section[\texorpdfstring{\bn{উপন্যাস} (upônyās)}{উপন্যাস (upônyās)}]{\bn{উপন্যাস} (upônyās) — a novel}

\label{lesson:BN-C327-uponyas}



\begin{quote}
Before the new one: say the Bengali for a hidden treasure, then the Bengali for to vanish.
\end{quote}

\subsection*{You'll want to know: \bn{উপন্যাস}}
\textbf{\bn{উপন্যাস}} — \emph{upônyās} — \textquotedblleft{}a novel\textquotedblright{}.

\begin{quote}
\bn{আমি} \bn{ছুটিতে} \bn{একটা} \bn{লম্বা} \bn{উপন্যাস} \bn{পড়লাম।}
\end{quote}

\begin{quote}
\emph{āmi chhuṭite ekṭā lômbā upônyās pôṛlām.}
\end{quote}

\begin{quote}
I read a long novel during the holidays.
\end{quote}

\begin{cousinweb}
\textbf{\bn{উপন্যাস}} (\emph{upônyās}) is Sanskrit \emph{upanyāsa}, 'a putting down, a placing near' (Monier-Williams); the sense 'a novel' is the modern Bengali one.
\end{cousinweb}

\begin{grammarlens}[title={What you've built}]
One more word for telling what happened.
\end{grammarlens}

\subsection*{Your turn}
\begin{itemize}
\raggedright
  \item \emph{Say it:} \emph{upônyās}
  \item \emph{Say it:} \emph{upônyās}, once more
  \item \emph{Say it:} \emph{āmi chhuṭite ekṭā lômbā upônyās pôṛlām.} — I read a long novel during the holidays.
  \item \emph{Recall:} say the Bengali for a hidden treasure, then the Bengali for to vanish, then say \emph{upônyās} again
\end{itemize}

\begin{tcolorbox}[breakable,colback=teal!4,colframe=teal!35!black,title={Before you move on}]
What does \bn{উপন্যাস} mean? (\textquotedblleft{}A novel\textquotedblright{}.) Say it once more.
\end{tcolorbox}
\section[\texorpdfstring{\bn{জীবনী} (jībônī)}{জীবনী (jībônī)}]{\bn{জীবনী} (jībônī) — a biography}

\label{lesson:BN-C327-jiboni}



\begin{quote}
Before the new one: say the Bengali for to vanish, then the Bengali for a novel.
\end{quote}

\subsection*{You'll want to know: \bn{জীবনী}}
\textbf{\bn{জীবনী}} — \emph{jībônī} — \textquotedblleft{}a biography\textquotedblright{}.

\begin{quote}
\bn{তিনি} \bn{এক} \bn{বিখ্যাত} \bn{বিজ্ঞানীর} \bn{জীবনী} \bn{লিখলেন।}
\end{quote}

\begin{quote}
\emph{tini ek bikhyātô bigyānīr jībônī likhlen.}
\end{quote}

\begin{quote}
He wrote a biography of a famous scientist.
\end{quote}

\begin{cousinweb}
\textbf{\bn{জীবনী}} (\emph{jībônī}) is \textbf{\bn{জীবন}}, life — Sanskrit \emph{jīvana} (Monier-Williams) — with the ending -ī: the written story of a life.
\end{cousinweb}

\begin{grammarlens}[title={What you've built}]
One more word for telling what happened.
\end{grammarlens}

\subsection*{Your turn}
\begin{itemize}
\raggedright
  \item \emph{Say it:} \emph{jībônī}
  \item \emph{Say it:} \emph{jībônī}, once more
  \item \emph{Say it:} \emph{tini ek bikhyātô bigyānīr jībônī likhlen.} — He wrote a biography of a famous scientist.
  \item \emph{Recall:} say the Bengali for to vanish, then the Bengali for a novel, then say \emph{jībônī} again
\end{itemize}

\begin{tcolorbox}[breakable,colback=teal!4,colframe=teal!35!black,title={Before you move on}]
What does \bn{জীবনী} mean? (\textquotedblleft{}A biography\textquotedblright{}.) Say it once more.
\end{tcolorbox}
\section[\texorpdfstring{\bn{গ্রন্থ} (grônthô)}{গ্রন্থ (grônthô)}]{\bn{গ্রন্থ} (grônthô) — a volume; a book (formal)}

\label{lesson:BN-C327-grontho}



\begin{quote}
Before the new one: say the Bengali for a novel, then the Bengali for a biography.
\end{quote}

\subsection*{You'll want to know: \bn{গ্রন্থ}}
\textbf{\bn{গ্রন্থ}} — \emph{grônthô} — \textquotedblleft{}a volume; a book (formal)\textquotedblright{}.

\begin{quote}
\bn{লাইব্রেরিতে} \bn{কবির} \bn{পুরনো} \bn{গ্রন্থগুলো} \bn{আছে।}
\end{quote}

\begin{quote}
\emph{lāibrerite kôbir purôno grônthôgulo āchhe.}
\end{quote}

\begin{quote}
The library has the poet's old volumes.
\end{quote}

\begin{cousinweb}
\textbf{\bn{গ্রন্থ}} (\emph{grônthô}) is Sanskrit \emph{grantha}, 'tying, binding, stringing together' (Monier-Williams): leaves strung together into a book. It is the formal word; \textbf{\bn{বই}} is the everyday one.
\end{cousinweb}

\begin{grammarlens}[title={What you've built}]
One more word for telling what happened.
\end{grammarlens}

\subsection*{Your turn}
\begin{itemize}
\raggedright
  \item \emph{Say it:} \emph{grônthô}
  \item \emph{Say it:} \emph{grônthô}, once more
  \item \emph{Say it:} \emph{lāibrerite kôbir purôno grônthôgulo āchhe.} — The library has the poet's old volumes.
  \item \emph{Recall:} say the Bengali for a novel, then the Bengali for a biography, then say \emph{grônthô} again
\end{itemize}

\begin{tcolorbox}[breakable,colback=teal!4,colframe=teal!35!black,title={Before you move on}]
What does \bn{গ্রন্থ} mean? (\textquotedblleft{}A volume; a book (formal)\textquotedblright{}.) Say it once more.
\end{tcolorbox}
\section[\texorpdfstring{\bn{অধ্যায়} (ôdhyāy)}{অধ্যায় (ôdhyāy)}]{\bn{অধ্যায়} (ôdhyāy) — a chapter of a book; a section}

\label{lesson:BN-C327-odhyay}



\begin{quote}
Before the new one: say the Bengali for a biography, then the Bengali for a volume.
\end{quote}

\subsection*{You'll want to know: \bn{অধ্যায়}}
\textbf{\bn{অধ্যায়}} — \emph{ôdhyāy} — \textquotedblleft{}a chapter of a book; a section\textquotedblright{}.

\begin{quote}
\bn{উপন্যাসের} \bn{প্রথম} \bn{অধ্যায়ে} \bn{নায়ক} \bn{গ্রাম} \bn{ছেড়ে} \bn{শহরে} \bn{আসে।}
\end{quote}

\begin{quote}
\emph{upônyāser prothom ôdhyāye nāyôk grām chheṛe shôhôre āse.}
\end{quote}

\begin{quote}
In the novel's opening section the hero leaves his village for the city.
\end{quote}

\begin{cousinweb}
\textbf{\bn{অধ্যায়}} (\emph{ôdhyāy}) is Sanskrit \emph{adhyāya}, 'a lesson, a lecture; a reading' (Monier-Williams), a part of a text read at one sitting.
\end{cousinweb}

\begin{grammarlens}[title={What you've built}]
One more word for telling what happened.
\end{grammarlens}

\subsection*{Your turn}
\begin{itemize}
\raggedright
  \item \emph{Say it:} \emph{ôdhyāy}
  \item \emph{Say it:} \emph{ôdhyāy}, once more
  \item \emph{Say it:} \emph{upônyāser prothom ôdhyāye nāyôk grām chheṛe shôhôre āse.} — In the novel's opening section the hero leaves his village for the city.
  \item \emph{Recall:} say the Bengali for a biography, then the Bengali for a volume, then say \emph{ôdhyāy} again
\end{itemize}

\begin{tcolorbox}[breakable,colback=teal!4,colframe=teal!35!black,title={Before you move on}]
What does \bn{অধ্যায়} mean? (\textquotedblleft{}A chapter of a book; a section\textquotedblright{}.) Say it once more.
\end{tcolorbox}
\section[\texorpdfstring{\bn{প্রকাশক} (prôkāshôk)}{প্রকাশক (prôkāshôk)}]{\bn{প্রকাশক} (prôkāshôk) — a publisher}

\label{lesson:BN-C327-prokashok}



\begin{quote}
Before the new one: say the Bengali for a volume, then the Bengali for a chapter of a book.
\end{quote}

\subsection*{You'll want to know: \bn{প্রকাশক}}
\textbf{\bn{প্রকাশক}} — \emph{prôkāshôk} — \textquotedblleft{}a publisher\textquotedblright{}.

\begin{quote}
\bn{প্রকাশক} \bn{লেখকের} \bn{নতুন} \bn{উপন্যাস} \bn{ছাপলেন।}
\end{quote}

\begin{quote}
\emph{prôkāshôk lekhôker nôtun upônyās chhāplen.}
\end{quote}

\begin{quote}
The publisher printed the writer's new novel.
\end{quote}

\begin{cousinweb}
\textbf{\bn{প্রকাশক}} (\emph{prôkāshôk}) is Sanskrit \emph{prakāśaka}, 'clear, shining; well known' (Monier-Williams), the agent noun beside \textbf{\bn{প্রকাশ} \bn{করা}}; the sense 'a publisher' is the modern one.
\end{cousinweb}

\begin{grammarlens}[title={What you've built}]
One more word for telling what happened.
\end{grammarlens}

\subsection*{Your turn}
\begin{itemize}
\raggedright
  \item \emph{Say it:} \emph{prôkāshôk}
  \item \emph{Say it:} \emph{prôkāshôk}, once more
  \item \emph{Say it:} \emph{prôkāshôk lekhôker nôtun upônyās chhāplen.} — The publisher printed the writer's new novel.
  \item \emph{Recall:} say the Bengali for a volume, then the Bengali for a chapter of a book, then say \emph{prôkāshôk} again
\end{itemize}

\begin{tcolorbox}[breakable,colback=teal!4,colframe=teal!35!black,title={Before you move on}]
What does \bn{প্রকাশক} mean? (\textquotedblleft{}A publisher\textquotedblright{}.) Say it once more.
\end{tcolorbox}
